\documentclass[10pt]{article}
\usepackage[T1]{fontenc}
\usepackage[utf8]{inputenc}
\usepackage[french]{babel}
\usepackage{amsmath,amssymb,mathrsfs,mathpazo}
\usepackage{geometry}
\geometry{a4paper,margin=2cm}
\usepackage{graphicx,array,tabularx,tabularray,colortbl,multirow,multicol,booktabs,enumitem,pifont,ulem,fancybox,diagbox,cancel}
\usepackage{pstricks,pstricks-add,pst-plot,pst-tree,pst-node,pst-text,pst-func,pst-math,pst-bspline,pst-3dplot,pst-eucl,pst-solides3d}
\let\pscalculate\undefined
\usepackage{pst-bezier}
\usepackage{tikz,pgfplots,tkz-tab}
\usetikzlibrary{arrows,arrows.meta,patterns,calc,positioning,trees,decorations.markings,shapes.geometric,shapes.misc,shapes.arrows,intersections,decorations.pathmorphing,decorations.pathreplacing,backgrounds,shadows,babel}
\pgfplotsset{compat=1.18}
\usepgfplotslibrary{fillbetween}
\usepackage[np]{numprint}
\usepackage{listings,makecell,pas-tableur,twoopt}
\let\barre\undefined
\usepackage[pstricks]{tablvar}
\definecolor{darkgreen}{RGB}{0,100,0}
\providecommand{\red}{\color{red}}
\providecommand{\blue}{\color{blue}}
\providecommand{\green}{\color{green}}
\providecommand{\black}{\color{black}}
\providecommand{\white}{\color{white}}
\providecommand{\gray}{\color{gray}}
\providecommand{\cyan}{\color{cyan}}
\providecommand{\magenta}{\color{magenta}}
\providecommand{\yellow}{\color{yellow}}
\providecommand{\violet}{\color{violet}}
\providecommand{\orange}{\color{orange}}
\providecommand{\euro}{\texteuro}
\providecommand{\EUR}{\texteuro}
\providecommand{\ieme}{\textsuperscript{e}}
\providecommand{\degres}{\ensuremath{^\circ}}
\providecommand{\degree}{\ensuremath{^\circ}}
\providecommand{\vv}[1]{\overrightarrow{#1}}
\providecommand{\Coord}[1]{#1}
\providecommand{\decosix}{}
\providecommand{\starredbullet}{\ensuremath{\bullet}}
\pagestyle{empty}
\setlength{\parindent}{0pt}
\usepackage{adjustbox}\DeclareUnicodeCharacter{2236}{:}\DeclareUnicodeCharacter{2190}{\ensuremath{\leftarrow}}
\providecommand{\R}{\mathbb{R}}
\providecommand{\N}{\mathbb{N}}
\providecommand{\Z}{\mathbb{Z}}
\providecommand{\Q}{\mathbb{Q}}
\providecommand{\D}{\mathbb{D}}
\providecommand{\C}{\mathbb{C}}
\providecommand{\e}{{\mathrm{e}}}
\providecommand{\vect}[1]{\overrightarrow{#1}}
\providecommand{\vectt}[1]{\overrightarrow{\mathrm{#1}}}
\providecommand{\vvt}[1]{\overrightarrow{\mathrm{#1}}}
\providecommand{\np}[2][]{{#2~#1}}
\providecommand{\vv}[1]{{\overrightarrow{#1}}}
\renewcommand{\d}{{\mathrm{d}}}
\providecommand{\degres}{{^\circ}}
\providecommand{\degre}{{^\circ}}
\providecommand{\blue}{\color{blue}}
\providecommand{\red}{\color{red}}
\providecommand{\enskip}{\quad}
\providecommand{\llbracket}{\lbrack\!\lbrack}
\providecommand{\rrbracket}{\rbrack\!\rbrack}
\providecommand{\euro}{{\text{€}}}
\providecommand{\no}{{\text{n°}}}
\providecommand{\interval}[2]{{[#1;#2]}}
\providecommand{\intervFO}[2]{{[#1;#2[}}
\providecommand{\dotfill}{\ldots}

\providecommand{\ds}{\displaystyle}
\providecommand{\textding}[1]{#1}
\providecommand{\Oij}{\(\mathrm{O};\vec{\imath},\vec{\jmath}\)}
\providecommand{\Oijk}{\(\mathrm{O};\vec{\imath},\vec{\jmath},\vec{k}\)}
\providecommand{\Ouv}{\(\mathrm{O};\vec{u},\vec{v}\)}
\providecommand{\intervFF}[2]{\ensuremath{[#1;#2]}}
\renewcommand{\euro}{\ifmmode\text{\texteuro}\else\texteuro\fi}
\DeclareUnicodeCharacter{00B0}{\ifmmode^{\circ}\else\textdegree\fi}
\DeclareUnicodeCharacter{20AC}{\ifmmode\text{\texteuro}\else\texteuro\fi}
% Énoncés : Éducation nationale. Transcription et corrigés : APMEP, Denis Vergès et contributeurs. Découpage et classement : CoopMaths, Sébastien Lozano, Éric Elter et contributeurs. Adaptation de la présentation : Maths Explorer.
% CC BY-SA : https://coopmaths.fr/www/about/
% https://coopmaths.fr/alea/static/bac/2026/tex/bac_2026_06_sujet1_etrangers_4.tex
\begin{document}
\noindent\textit{APMEP / CoopMaths. CC BY-SA. Présentation adaptée par Maths Explorer.}\par\medskip
On considère la fonction $f$ définie sur l'intervalle $] 0 ; +\infty[$ par \quad \[f(x)=\dfrac{\ln (x)}{x^{2}}\]
	
	On note $\mathcal{C}_{f}$ sa courbe représentative dans un repère orthonormé du plan.
	
	On admet que la fonction $f$ est deux fois dérivable sur l'intervalle $]0 ; +\infty[$.
	
	On note $f'$ sa fonction dérivée et $f''$ sa fonction dérivée seconde.
	
	\subsection*{Partie A}
	\begin{enumerate}
		\item \begin{enumerate}
			\item Déterminer les limites de la fonction $f$ aux bornes de l'intervalle $] 0;+\infty[$.
			
			\item Interpréter graphiquement les résultats obtenus.
		\end{enumerate}
		
		\item Pour tout réel strictement positif $x$, démontrer que \quad
		\[f'(x)=\dfrac{1-2 \ln (x)}{x^{3}}\]
		
		\item Dresser le tableau de variation de la fonction $f$ en précisant les limites et les valeurs exactes des éventuels extremums de la fonction.
		
		\item Déterminer l'équation réduite de la droite $\Delta$, tangente à la courbe $\mathcal{C}_{f}$ au point A d'abscisse 1.
		
		\item Vérifier que pour tout réel strictement positif $x$, on a \quad \[f''(x)=\dfrac{-5+6 \ln (x)}{x^{4}}\]
		
		\item \begin{enumerate}
			\item Étudier la convexité de la fonction $f$ en précisant les coordonnées des éventuels points d'inflexion.
			
			\item En déduire que, pour tout réel $x$ appartenant à $\left]0 ; \e^{\frac{5}{6}}\right]$, on a \quad
			\[x-1 \geqslant \dfrac{\ln (x)}{x^{2}}\]
		\end{enumerate}
		
		\item Justifier que, pour tout réel $x$ appartenant à $\left[\e^{\frac{5}{6}} ; +\infty\right[$, on a aussi
		\[x-1 \geqslant  \dfrac{\ln (x)}{x^{2}}\]
	\end{enumerate}
	
	\subsection*{Partie B}
	On considère la suite $(I_{n})$ définie, pour tout entier naturel non nul $n$, par :
	\[I_{n}=\int_{1}^{n} \dfrac{\ln (x)}{x^{2}} \d x\]
	
	\begin{enumerate}
		\item Donner une interprétation graphique de $I_{n}$ pour $n$ un entier naturel non nul.
		
		\item Démontrer que la suite $(I_{n})$ est croissante.
		
		\item On souhaite calculer une approximation de \quad $\ds \int_{1}^{10} f(x) \d x$ \quad en déterminant la somme des aires des rectangles tracés dans le graphique ci-dessous.
		
		\begin{center}
			\begin{tikzpicture}
				\begin{axis}[/pgf/number format/.cd, use comma,
					x={10mm}, y={20cm},
					xmin = -0.9, xmax = 11.9,
					ymin = -0.04, ymax = 0.28,
					xtick distance=1,
					ytick distance=0.05,
					scaled ticks=false,
					tick label style={/pgf/number format/fixed, /pgf/number format/precision=2,font=\footnotesize},
					minor tick num = 0,
					grid=major,axis lines=center]
					\node[below left=2pt] at (axis cs: 0,0) {\footnotesize 0};
					\node[red] at(axis cs: 10.5,0.035) {$ \mathcal{C}_f $};
					\pgfplotsextra{
						\foreach \x/\y in {1/0.184, 2/0.173, 3/0.122, 4/0.087, 5/0.064, 6/0.05, 7/0.04, 8/0.032, 9/0.027}
						\draw[brown,thick,fill opacity=0.2,fill] (axis cs: \x,\y) rectangle (axis cs: {\x+1},0) ;
					}
					\draw[thick,blue,shift={(axis cs: 1.649,0.184)}] (3pt,3pt)--(-3pt,-3pt) (3pt,-3pt)--(-3pt,3pt) ;
					\addplot [line width=1.2pt,color=teal,smooth,samples=100,domain= 0.8:12 ]{ln(x)/(x*x)};
				\end{axis}
			\end{tikzpicture}
		\end{center}
		
		
		Pour cela, on utilise le script suivant, écrit en langage Python.
		
		\begin{center}
			\begin{tblr}{colspec={X[6cm,font=\ttfamily]},hline{1,Z}={solid}, vlines}
				from math import *\\
				\\
				S=1/(2*exp(1))\\
				for i in range (2,10):\\
				\quad S =...\\
				print(S)\\
			\end{tblr}
		\end{center}
		
		Recopier et compléter ce script afin qu'il réponde au problème posé.
		
		\item À l'aide d'une intégration par parties, démontrer que pour tout entier naturel non nul $n$,
		\[I_{n}=\dfrac{n-1-\ln (n)}{n}\]
		
		\item Calculer la limite de la suite $(I_{n})$ quand $n$ tend vers $+\infty$.
	\end{enumerate}

\end{document}
